% Einheit 1 von 3 – Schnittpunkt von Gerade und Ebene (bank/schnittmengen/e1.jsonl, zusammenbau v0.1)
%% TODO Zeitmarke Einheit 1: Marken-Zeile nicht lesbar
%% TODO Zweigzeile Teil 1: was hier gelernt wird (Satz in Schülersprache aus dem Einheitstitel) – Einheit 1
\einheitenkopf[e1]{Einheit 1 von 3 · Schnittpunkt von Gerade und Ebene}
\zweigzeile{Abitur GK}
\setcounter{aufgabe}{8}

%% TODO Titel: Ich-kann-Satz fehlt in der Bank (Platzhalter: Kettenname) – Nr. 9
%% TODO Anweisung: ein Satz über der ersten Teilaufgabe fehlt in der Bank – Nr. 9
\begin{aufgabe}{Was wird geschnitten – und was kann herauskommen?}
\begin{teile}
\teil Eine Gerade durchstößt eine Ebene, sie ist nicht parallel zu ihr. Was ist die Schnittmenge? \\ \kreuz{ein Punkt} \\ \kreuz{eine Gerade} \\ \kreuz{leer}
\end{teile}
\end{aufgabe}

%% TODO Titel: Ich-kann-Satz fehlt in der Bank (Platzhalter: Kettenname) – Nr. 10
%% TODO Anweisung: ein Satz über der ersten Teilaufgabe fehlt in der Bank – Nr. 10
\begin{aufgabe}{Schnittpunkt von Gerade und Ebene}
%% TODO \rechenplatz{n}: Schrittzahl des Grundfalls fehlt in der Bank (nur bei fester Schreibform, 2.3 b)
\begin{teile}
\teil Gegeben sind eine Gerade in Parameterform und eine Ebene in Koordinatenform. Welcher Weg führt zum Schnittpunkt? \\ \kreuz{einsetzen} \\ \kreuz{gleichsetzen}
\teil Bestimme den Schnittpunkt der Geraden $\vec{x} = (-1|2|1) + t \cdot (1|0|1)$ mit der Ebene $x + y + z = 6$. S( \leerfeld | \leerfeld | \leerfeld )
\teil Ein Körper hat die Spitze $P(0|0|3)$; $M(3|1|0)$ ist der Mittelpunkt einer Grundkante. Die Gerade durch $P$ und $M$ schneidet die Ebene $L$: $x + 2y + z = 4$ im Punkt $Q$. Bestimme die Koordinaten von $Q$. \hfill (Abitur 2023 GK) \\ Q( \leerfeld | \leerfeld | \leerfeld )
\teil Ein Ausstellungsgebäude wird durch einen Körper modelliert; das Dreieck $DEF$ mit $D(1|0|6)$, $E(1|8|6)$ und $F(-5|2|6)$ liegt parallel zur $xy$-Ebene. Weise nach, dass sich die Gerade durch $A(-2|2|0)$ und $G(-4|4|8)$ und die Ebene, in der das Dreieck $DEF$ liegt, im Punkt $R(-3{,}5|3{,}5|6)$ schneiden. \hfill (Abitur 2018 LK)
\teil Ein Sonnensegel hat die Ecken $A(1|2|2)$, $B(1|6|2)$ und $C(-2|4|4)$; der Boden ist die $xy$-Ebene, eine Längeneinheit entspricht 1 m. Bei parallelem Sonnenlicht liegen die Schatten von $A$ und $B$ in $A'(-1|3|0)$ und $B'(-1|7|0)$. Bestimme den Schatten $C'$ von $C$. \hfill (Abitur 2020 GK) \\ C'( \leerfeld | \leerfeld | \leerfeld )
\teil Ein Rasenroboter fährt geradlinig und gleichförmig von $P(8|2|0)$ über $Q(6|3|0)$ weiter; für die Strecke von $P$ nach $Q$ braucht er 4 s. Eine Lichtschranke liegt in der senkrechten Ebene $E$: $x - 2y = -6$. Wie lange braucht der Roboter von $Q$ bis zur Lichtschranke? \hfill (Abitur 2018 GK) \\ \leerfeld[s]
\teil Eine Mauer hat die obere Kante $PQ$ mit $P(6|-2|2)$ und $Q(6|8|2)$. Ein Beobachter hat das Auge im Punkt $K(9|3|1{,}5)$. Hinter der Mauer steht ein Gerätehaus; der Teil seiner Seitenwand, den der Beobachter nicht sehen kann, ist das Dreieck $GBH$, wobei $H$ auf der Kante $BE$ liegt. Beschreibe, wie die Koordinaten von $H$ rechnerisch bestimmt werden können. \hfill (Abitur 2023 LK)
\end{teile}
\end{aufgabe}

%% TODO Titel: Ich-kann-Satz fehlt in der Bank (Platzhalter: Kettenname) – Nr. 11
%% TODO Anweisung: ein Satz über der ersten Teilaufgabe fehlt in der Bank – Nr. 11
\begin{aufgabe}{Höhe einer Spitze aus dem Abstand ihres Schattenpunkts zu einem Punkt bestimmen}
\begin{teile}
\teil Ein Fahnenmast steht senkrecht im Ursprung, seine Spitze ist $S(0|0|h)$ mit $h > 0$; der Boden ist die $xy$-Ebene, eine Längeneinheit entspricht 1 m. Sonnenlicht fällt in parallelen Geraden mit dem Richtungsvektor $(4|2|-4)$. Der Schatten der Spitze liegt auf dem Boden und hat vom Punkt $B(1|-2|0)$ den Abstand 5 m. Bestimme die Höhe $h$. \hfill (Abitur 2018 LK) \\ h = \leerfeld[m]
\end{teile}
\end{aufgabe}

%% TODO Titel: Ich-kann-Satz fehlt in der Bank (Platzhalter: Kettenname) – Nr. 12
%% TODO Anweisung: ein Satz über der ersten Teilaufgabe fehlt in der Bank – Nr. 12
\begin{aufgabe}{Schnittpunkt von Gerade und Ebene – Fehler finden, Begründen}
\begin{teile}
\teil Ich finde den Fehler: Gesucht ist der Schnittpunkt von $\vec{x} = (1|2|0) + t \cdot (1|1|2)$ mit der Ebene $x + y + z = 7$. Ole setzt ein: \rechnung{t + t + 2t &= 7 \\ t &= 1{,}75}
\teil Begründe, warum das Einsetzen des allgemeinen Geradenpunkts in eine Koordinatengleichung eine Gleichung mit nur einer Unbekannten liefert.
\end{teile}
\end{aufgabe}
